\documentclass[12pt,a4paper]{article}
\usepackage[spanish]{babel}
\usepackage[utf8]{inputenc}
\usepackage{amsmath, amssymb}
\usepackage{geometry}
\usepackage{multicol}
\usepackage{fancyhdr}
\usepackage{graphicx}

\geometry{margin=2cm}
\pagestyle{fancy}
\fancyhf{}
\rhead{Máquinas Eléctricas II}
\lhead{Evaluación - 6°1°}
\rfoot{Colegio Secundario Técnico}

\begin{document}

\begin{center}
{\Large \textbf{Evaluación de Máquinas Eléctricas II}}\\[1em]
\textbf{Alumno:} \underline{\hspace{7cm}} \hfill \textbf{Curso:} 6°1° \\[0.5em]
\textbf{Fecha:} \underline{\hspace{4cm}} \hfill \textbf{Profesor:} Gustavo \\[1em]
\hrule
\end{center}

\vspace{1em}

\begin{center}
{\Large \textbf{Tema 1: Alternador Trifásico (TP N°7)}}\\[0.5em]
\textit{Contenido: Alternadores, frecuencia, tensiones simples y compuestas, cargas trifásicas.}
\end{center}

\begin{enumerate}
\item Determinar la frecuencia que produce un alternador que gira a \( n = 1500\ \text{rpm} \) y posee dos pares de polos.
\item Calcular la tensión compuesta de un sistema trifásico con tensión simple de \( V_f = 223\ \text{V} \).
\item Un motor trifásico conectado en estrella desarrolla una potencia de \( P = 10\,kW \), con \( \cos(\varphi) = 0{,}8 \) y \( V_L = 400\ \text{V} \). Hallar las corrientes de línea y fase.
\item Se conectan en estrella tres bobinas iguales con \( R = 12\,\Omega \) y \( X_L = 20\,\Omega \), alimentadas con \( V_L = 220\,V \). Calcular:
\[
Z,\quad I_f,\quad \cos(\varphi),\quad P,\quad Q,\quad S
\]
\item Un calefactor trifásico consta de tres resistencias de \( R = 10\,\Omega \) conectadas:
\begin{enumerate}
  \item En estrella.
  \item En triángulo.
\end{enumerate}
Determinar \( I_f \), \( I_L \) y la potencia consumida en cada caso.
\end{enumerate}

\vspace{1em}
\hrule
\vspace{1em}

\begin{center}
{\Large \textbf{Tema 2: Sistema Trifásico Desbalanceado (TP N°8)}}\\[0.5em]
\textit{Cálculo de corriente de neutro.}
\end{center}

\begin{enumerate}
\item Explicar qué es un sistema trifásico desbalanceado y por qué aparece la corriente de neutro.
\item En un sistema trifásico con \( I_A = 12\,A \), \( I_B = 9\,A \), \( I_C = 5\,A \) y desfase \( \theta = 120^\circ \), calcular la corriente de neutro usando:
\[
I_N = \sqrt{I_A^2 + I_B^2 + I_C^2 + 2(I_A I_B + I_B I_C + I_C I_A)\cos(120^\circ)}
\]
\item Si \( I_A = 20\,A \), \( I_B = 18\,A \), \( I_C = 15\,A \), determinar \( I_N \).
\item ¿Qué consecuencias técnicas puede tener una corriente de neutro elevada en una instalación trifásica?
\item Resolver con números complejos: 
\[
V_A = 230\angle0^\circ,\quad V_B = 225\angle-120^\circ,\quad V_C = 220\angle120^\circ
\]
\[
Z_A = 16+j9,\quad Z_B = 24+j6,\quad Z_C = 18+j9
\]
Hallar las corrientes de fase y la corriente de neutro.
\end{enumerate}

\vspace{2em}
\begin{flushright}
\textbf{--- Fin de la evaluación ---}
\end{flushright}

\end{document}
